\chapter{Komutant operatora i centralizatory macierzy}
\label{ch:lin21-commutant}
\index{komutant operatora@Komutant operatora}
\index{centralizator macierzy@Centralizator macierzy}
\index{endomorfizm modułu@Endomorfizm modułu}
\index{centralizator odwracalny@Centralizator odwracalny}

\noindent\textbf{Cel rozdziału.}
W Rozdziałach~\ref{ch:lin19-frobenius}
i~\ref{ch:lin20-elementary}
sklasyfikowaliśmy endomorfizmy
z dokładnością do podobieństwa.
Teraz badamy wszystkie operatory, które
\emph{przemieniają się} z ustalonym operatorem.
Pokażemy, że ich przestrzeń jest
algebrą endomorfizmów modułu nad \(K[t]\).
Obliczymy jej wymiar za pomocą
czynników niezmienniczych, dzielników
elementarnych i wymiarów jąder wielomianów
operatora. Przedstawimy także przykłady
centralizatorów odwracalnych.

W całym rozdziale \(K\) jest dowolnym
ciałem, \(V\) jest skończenie wymiarową
przestrzenią wektorową nad \(K\),
\(n=\dim_K V\) i
\(T\in\operatorname{End}_K(V)\).
Dopuszczamy \(n=0\): wtedy
\(\operatorname{End}_K(V)\) jest zerową
przestrzenią wektorową, a jedyny
endomorfizm jest zarazem identycznością
i odwzorowaniem zerowym.

\section{Komutant i jego podstawowe własności}
\label{sec:lin21-definition}

\begin{definicja}[Komutant i centralizator odwracalny]
\label{def:lin21-commutant}
Dla \(T\in\operatorname{End}_K(V)\)
definiujemy jego \emph{komutant}
jako zbiór
\[
 \mathcal C(T)
 =\{S\in\operatorname{End}_K(V):ST=TS\}.
\]
\emph{Centralizatorem odwracalnym}
operatora \(T\) nazywamy
\[
 \mathcal Z(T)
 =\{S\in\operatorname{GL}_K(V):ST=TS\}
 =\mathcal C(T)\cap\operatorname{GL}_K(V).
\]
Dla macierzy \(A\in\operatorname{Mat}_{n\times n}(K)\)
stosujemy analogiczne oznaczenia
\(\mathcal C(A)\) i \(\mathcal Z(A)\).
\end{definicja}

\begin{stwierdzenie}[Algebra komutanta i niezmienniczość]
\label{prop:lin21-algebra}
Dla każdego \(T\in\operatorname{End}_K(V)\)
zbiór \(\mathcal C(T)\) jest
podalgebrą \(\operatorname{End}_K(V)\),
zawierającą \(I_V\) i wszystkie
wielomiany \(p(T)\), \(p\in K[t]\).
Zbiór \(\mathcal Z(T)\) jest grupą
elementów odwracalnych tej algebry.

Jeżeli \(U\colon V\to W\)
jest izomorfizmem \(K\)-liniowym
i \(T'=UTU^{-1}\), to
\[
 \mathcal C(T')=U\mathcal C(T)U^{-1},
 \qquad
 \mathcal Z(T')=U\mathcal Z(T)U^{-1}.
\]
W szczególności wymiar komutanta
zależy jedynie od klasy podobieństwa
operatora.
\end{stwierdzenie}
\begin{proof}
Jeśli \(ST=TS\) i \(RT=TR\),
to dla dowolnych \(a,b\in K\)
mamy \((aS+bR)T=T(aS+bR)\),
a także
\((SR)T=S(TR)=(ST)R=T(SR)\).
Ponadto \(I_VT=TI_V\).
Indukcja po stopniu pokazuje
\(T^jT=TT^j\), skąd
\(p(T)T=Tp(T)\).
Otrzymujemy podalgebrę z jedynką.

Jeżeli \(S\in\mathcal C(T)\)
jest odwracalny, to z \(ST=TS\)
przez pomnożenie z lewej i prawej
przez \(S^{-1}\) otrzymujemy
\(S^{-1}T=TS^{-1}\).
Elementy odwracalne komutanta
tworzą więc grupę \(\mathcal Z(T)\).

Dla \(S'=USU^{-1}\) zachodzi
\[
 S'T'=USTU^{-1},\qquad
 T'S'=UTSU^{-1}.
\]
Zatem \(S'T'=T'S'\) dokładnie wtedy,
gdy \(ST=TS\).
Sprzężenie przez \(U\) zachowuje
również odwracalność i jest
izomorfizmem \(K\)-liniowym.
To daje obie równości i tezę o wymiarze.
\end{proof}

\begin{przyklad}[Macierz skalarna i macierz diagonalna]
\label{prz:lin21-scalar}
Dla \(A=\lambda I_n\) każdy
\(X\in\operatorname{Mat}_{n\times n}(K)\)
spełnia \(XA=AX\).
Dlatego
\[
 \mathcal C(\lambda I_n)
 =\operatorname{Mat}_{n\times n}(K),
 \qquad \dim_K\mathcal C(\lambda I_n)=n^2,
 \qquad \mathcal Z(\lambda I_n)=\operatorname{GL}_n(K).
\]
Jeśli \(A=\operatorname{diag}(\lambda_1,\ldots,\lambda_n)\),
to dla \(X=(x_{ij})\) warunek
\(XA=AX\) jest równoważny
\[
 (\lambda_j-\lambda_i)x_{ij}=0
 \quad\text{dla wszystkich }i,j.
\]
Zatem \(x_{ij}\) może być różny od zera
tylko wtedy, gdy \(\lambda_i=\lambda_j\).
Jeżeli różne wartości diagonalne
\(\lambda\) występują z krotnościami
\(m_\lambda\), komutant składa się
z bloków \(\operatorname{Mat}_{m_\lambda}(K)\),
a jego wymiar wynosi
\(\sum_\lambda m_\lambda^2\).
\end{przyklad}

\section{Odwzorowania liniowe między modułami cyklicznymi}
\label{sec:lin21-cyclic}

\begin{stwierdzenie}[Komutant a endomorfizmy modułu]
\label{prop:lin21-module-endo}
Jeżeli \(V_T\) jest modułem
nad \(K[t]\) z
Definicji~\ref{def:lin19-module},
to
\[
 \mathcal C(T)
 =\operatorname{End}_{K[t]}(V_T)
\]
jako zbiory odwzorowań i jako algebry
nad \(K\), przy utożsamieniu
obu stron z odpowiednimi operatorami na \(V\).
\end{stwierdzenie}
\begin{proof}
Każdy homomorfizm modułów nad \(K[t]\)
jest \(K\)-liniowy, ponieważ
skalary \(K\) są wielomianami stałymi.
Zachowuje również działanie \(t\),
a więc spełnia \(ST=TS\).

Odwrotnie, jeśli \(S\) jest \(K\)-liniowy
i \(ST=TS\), to indukcja daje
\(ST^j=T^jS\) dla każdego \(j\ge0\).
Dla \(p(t)=\sum_j a_jt^j\)
wynika stąd
\(Sp(T)=p(T)S\).
A zatem \(S(p\cdot v)=p\cdot S(v)\)
dla wszystkich \(p\in K[t]\)
i \(v\in V\).
Zgodność dodawania i składania
jest bezpośrednia, więc otrzymujemy
równość algebr.
\end{proof}

\begin{twierdzenie}[Homomorfizmy modułów cyklicznych]
\label{thm:lin21-hom-cyclic}
Niech \(f,g\in K[t]\) będą monicznymi
wielomianami dodatnich stopni.
Niech \(d=\gcd(f,g)\) będzie
ich monicznym największym wspólnym
dzielnikiem.
Wówczas istnieje izomorfizm
przestrzeni wektorowych nad \(K\)
\[
 \operatorname{Hom}_{K[t]}
 \bigl(K[t]/(f),K[t]/(g)\bigr)
 \cong K[t]/(d).
\]
W szczególności
\[
 \dim_K\operatorname{Hom}_{K[t]}
 \bigl(K[t]/(f),K[t]/(g)\bigr)
 =\deg\gcd(f,g).
\]
\end{twierdzenie}
\begin{proof}
Niech \(\varphi\) będzie homomorfizmem
\(K[t]\)-modułów z lewej strony.
Każda klasa modulo \(f\) ma postać
\(h\cdot[1]_f\), więc
\(\varphi([h]_f)=h\varphi([1]_f)\).
Zatem \(\varphi\) jest całkowicie
wyznaczony przez obraz
\(y=\varphi([1]_f)\in K[t]/(g)\).

Ponieważ \(f\cdot[1]_f=0\),
konieczny jest warunek
\[
 f\cdot y=0\quad\text{w }K[t]/(g).
\]
Jest on również wystarczający:
dla takiego \(y\) wzór
\(\varphi_y([h]_f)=h\cdot y\)
nie zależy od wyboru przedstawiciela,
bo zmiana \(h\) o wielokrotność \(f\)
nie zmienia wyniku.
Otrzymujemy zatem izomorfizm
z przestrzenią
\(\{y\in K[t]/(g):fy=0\}\).

Zapiszmy \(f=df_0\), \(g=dg_0\).
Z definicji największego wspólnego
dzielnika wynika
\(\gcd(f_0,g_0)=1\).
Warunek \(g\mid fy\), dla dowolnego
wielomianowego przedstawiciela \(y\),
jest więc równoważny \(g_0\mid y\):
po skróceniu przez niezerowy \(d\)
otrzymujemy \(g_0\mid f_0y\),
a lemat Euklidesa daje \(g_0\mid y\).

Zdefiniujmy odwzorowanie
\[
 \Psi\colon K[t]/(d)\longrightarrow
 \{y\in K[t]/(g):fy=0\},\qquad
 [h]_d\longmapsto[g_0h]_g.
\]
Jeżeli \(h-h'\) jest wielokrotnością
\(d\), to \(g_0(h-h')\) jest
wielokrotnością \(g=dg_0\),
więc \(\Psi\) jest dobrze określone.
Jest surjektywne, bo każdy element
jądra mnożenia przez \(f\) ma
przedstawiciela podzielnego przez \(g_0\).
Jest różnowartościowe:
\([g_0h]_g=0\) oznacza
\(dg_0\mid g_0h\), a po skróceniu
\(d\mid h\).
Jest też \(K\)-liniowe.
Połączenie opisanych izomorfizmów
daje tezę; wymiar \(K[t]/(d)\)
wynosi \(\deg d\) na mocy
Stwierdzenia~\ref{prop:lin19-companion-module},
jeżeli \(\deg d>0\).
Gdy \(d=1\), obie przestrzenie
są zerowe i wzór na wymiar
również pozostaje poprawny.
\end{proof}

\begin{wniosek}[Komutant operatora cyklicznego]
\label{cor:lin21-cyclic}
Niech \(n=\dim_K V\ge1\) i
niech \(T\in\operatorname{End}_K(V)\)
ma wektor cykliczny.
Wówczas
\[
 \mathcal C(T)=K[T]
 =\{p(T):p\in K[t]\}
\]
oraz \(\dim_K\mathcal C(T)=n\).
Jeżeli \(m_T\) jest wielomianem
minimalnym, to odwzorowanie
\[
 K[t]/(m_T)\longrightarrow\mathcal C(T),
 \qquad [p]\longmapsto p(T)
\]
jest izomorfizmem algebr
z jedynką nad \(K\).
Ponadto \(p(T)\) jest odwracalny
dokładnie wtedy, gdy
\(\gcd(p,m_T)=1\).
\end{wniosek}
\begin{proof}
Z Wniosku~\ref{cor:lin19-cyclic}
wynika, że \(V_T\) jest
izomorficzny z \(K[t]/(m_T)\)
i \(\deg m_T=n\).
Niech \(S\in\mathcal C(T)\).
Wobec
Stwierdzenia~\ref{prop:lin21-module-endo}
odpowiada mu endomorfizm modułu
\(K[t]/(m_T)\).
Na podstawie dowodu
Twierdzenia~\ref{thm:lin21-hom-cyclic},
przy \(f=g=m_T\), jest on wyznaczony
przez dowolny obraz klasy \([1]\),
a więc jest mnożeniem przez pewną
klasę \([p]\).
Po przeniesieniu na \(V_T\)
jest to dokładnie operator \(p(T)\).
Zatem \(\mathcal C(T)=K[T]\).

Jądro homomorfizmu
\(K[t]\to\mathcal C(T)\),
\(p\mapsto p(T)\), jest ideałem
\((m_T)\) z definicji
wielomianu minimalnego.
Stąd indukowany homomorfizm
ilorazowy jest izomorfizmem algebr.
Wymiar jest równy
\(\deg m_T=n\).

Jeżeli \(\gcd(p,m_T)=1\),
to tożsamość Bézouta daje
\(up+vm_T=1\).
Po podstawieniu \(T\)
otrzymujemy \(u(T)p(T)=I_V\),
więc \(p(T)\) jest odwracalny.
Odwrotnie, jeśli \(p(T)\)
jest odwracalny, to jest
odwracalny także w komutancie
na mocy Stwierdzenia~\ref{prop:lin21-algebra}.
Jego odwrotność jest zatem
pewnym \(u(T)\), skąd
\((up-1)(T)=0\), czyli
\(m_T\mid up-1\).
Tożsamość \(up+vm_T=1\)
dla odpowiedniego \(v\)
wymusza \(\gcd(p,m_T)=1\).
\end{proof}

\begin{przyklad}[Pojedynczy blok Jordana]
\label{prz:lin21-jordan}
Niech \(J=J_3(\lambda)\) nad \(K\).
Ponieważ \(J=\lambda I_3+N\), gdzie
\(N^3=0\) i \(N^2\ne0\),
jego wielomian minimalny to
\((t-\lambda)^3\).
Pierwszy wektor bazy potęg
dla \(C((t-\lambda)^3)\)
jest cykliczny, zatem na mocy
Stwierdzenia~\ref{prop:lin20-linear-primary}
również \(J\) ma wektor cykliczny.
Wniosek~\ref{cor:lin21-cyclic}
daje więc
\[
 \mathcal C(J)
 =\{aI_3+bN+cN^2:a,b,c\in K\}.
\]
W zwykłej bazie Jordana oznacza to
dokładnie macierze
\[
 X=\begin{pmatrix}
 a&b&c\\0&a&b\\0&0&a
 \end{pmatrix}.
\]
Ponieważ \(X\) jest trójkątna,
\(\det X=a^3\), a zatem
\(X\in\mathcal Z(J)\) dokładnie
wtedy, gdy \(a\ne0\).
\end{przyklad}

\section{Wymiar komutanta z czynników niezmienniczych}
\label{sec:lin21-invariant}

\begin{twierdzenie}[Wzór na wymiar komutanta]
\label{thm:lin21-invariant-dimension}
Niech \(T\in\operatorname{End}_K(V)\),
a \(f_1\mid\cdots\mid f_r\)
niech będą jego czynnikami
niezmienniczymi z
Twierdzenia~\ref{thm:lin19-frobenius}.
Wówczas
\[
 \dim_K\mathcal C(T)
 =\sum_{i=1}^{r}\sum_{j=1}^{r}
       \deg\gcd(f_i,f_j)
 =\sum_{i=1}^{r}\bigl(2(r-i)+1\bigr)\deg f_i.
\]
Dla \(V=\{0\}\), gdy \(r=0\),
obie sumy są puste i mają wartość \(0\).
\end{twierdzenie}
\begin{proof}
Na mocy Twierdzenia~\ref{thm:lin19-frobenius}
istnieje izomorfizm modułów
\[
 V_T\cong\bigoplus_{i=1}^{r}M_i,
 \qquad M_i=K[t]/(f_i).
\]
Każdy endomorfizm tej skończonej
sumy prostej jest jednoznacznie
wyznaczony przez składowe
\(\varphi_{ij}\colon M_j\to M_i\),
otrzymane przez złożenie
włożenia \(M_j\) do sumy,
endomorfizmu i rzutu na \(M_i\).
Odwrotnie, dowolna rodzina
homomorfizmów modułów \(\varphi_{ij}\)
wyznacza endomorfizm sumy
wzorem
\((x_j)_j\mapsto(\sum_j\varphi_{ij}(x_j))_i\).
Stąd izomorfizm przestrzeni wektorowych
\[
 \operatorname{End}_{K[t]}(V_T)
 \cong\bigoplus_{i,j=1}^{r}
       \operatorname{Hom}_{K[t]}(M_j,M_i).
\]
Wykorzystując
Stwierdzenie~\ref{prop:lin21-module-endo}
i Twierdzenie~\ref{thm:lin21-hom-cyclic},
po zsumowaniu wymiarów
otrzymujemy pierwszy wzór.

Z podzielności \(f_1\mid\cdots\mid f_r\)
wynika
\(\gcd(f_i,f_j)=f_{\min(i,j)}\).
Dla każdego indeksu \(k\)
jest dokładnie \(2(r-k)+1\)
uporządkowanych par \((i,j)\),
dla których \(\min(i,j)=k\):
są to pary \((k,j)\) z \(j\ge k\)
oraz pary \((i,k)\) z \(i>k\).
Grupując składniki pierwszej sumy
według mniejszego indeksu,
dostajemy drugi wzór.
Przy \(r=0\) obie strony są zerowe.
\end{proof}

\begin{wniosek}[Oszacowanie wymiaru komutanta]
\label{cor:lin21-bounds}
Dla \(T\in\operatorname{End}_K(V)\),
gdzie \(n=\dim_K V\ge1\), zachodzi
\[
 n\le\dim_K\mathcal C(T)\le n^2.
\]
Równość po lewej zachodzi wtedy i tylko
wtedy, gdy \(T\) jest operatorem cyklicznym.
Równość po prawej zachodzi wtedy i tylko
wtedy, gdy \(T=\lambda I_V\)
dla pewnego \(\lambda\in K\).
\end{wniosek}
\begin{proof}
Niech \(f_1,\ldots,f_r\)
będą czynnikami niezmienniczymi.
Z Twierdzenia~\ref{thm:lin19-frobenius}
wynika \(n=\sum_i\deg f_i\),
a każdy \(\deg f_i\ge1\).
Wzór z
Twierdzenia~\ref{thm:lin21-invariant-dimension}
ma dodatnie współczynniki
\(2(r-i)+1\ge1\),
stąd \(\dim\mathcal C(T)\ge n\).
Jeżeli \(r=1\), otrzymujemy
równość. Jeżeli \(r>1\),
współczynnik przy \(\deg f_1\)
jest większy od \(1\),
więc nierówność jest ostra.
Równość po lewej jest zatem
równoważna \(r=1\), a to
równoważne cykliczności z
Wniosku~\ref{cor:lin19-cyclic}.

Komutant jest podprzestrzenią
\(n^2\)-wymiarowej przestrzeni
\(\operatorname{End}_K(V)\),
więc jego wymiar nie przekracza \(n^2\).
Jeśli jest równy \(n^2\),
to \(\mathcal C(T)=\operatorname{End}_K(V)\);
operator \(T\) przemienia się więc
z każdym endomorfizmem.
Wybierzmy bazę i oznaczmy
przez \(E_{ij}\) operator
odpowiadający jednostce macierzowej.
Z równości \(TE_{ii}=E_{ii}T\)
dla wszystkich \(i\)
wynika, że macierz \(T\)
jest diagonalna: dla \(k\ne i\)
jej współczynnik \((k,i)\)
oraz \((i,k)\) musi zniknąć.
Następnie dla \(i\ne j\)
równość \(TE_{ij}=E_{ij}T\)
daje \(t_{ii}=t_{jj}\).
Zatem wszystkie elementy diagonalne
są równe, czyli \(T=\lambda I_V\).
Dla \(n=1\) każdy operator
jest skalarny, więc wniosek
również zachodzi.
W drugą stronę operator skalarny
przemienia się z każdym endomorfizmem,
co pokazano także w
Przykładzie~\ref{prz:lin21-scalar}.
\end{proof}

\section{Dzielniki elementarne i filtracje jąder}
\label{sec:lin21-elementary}

\begin{twierdzenie}[Wzór z dzielników elementarnych]
\label{thm:lin21-elementary-dimension}
Niech dla każdego monicznego
nierozkładalnego \(p\in K[t]\),
występującego w dzielnikach
elementarnych operatora \(T\),
jego odpowiednie wykładniki
wynoszą \(a_{p,1},\ldots,a_{p,r_p}\ge1\).
Niech \(d_p=\deg p\)
oraz niech \(B_{p,k}\) oznacza
liczbę tych wykładników,
które są nie mniejsze od \(k\).
Wówczas
\[
 \dim_K\mathcal C(T)
 =\sum_p d_p\sum_{i,j=1}^{r_p}
                   \min(a_{p,i},a_{p,j})
 =\sum_p d_p\sum_{k\ge1} B_{p,k}^2.
\]
Wszystkie sumy są skończone.
Jeśli
\(D_{p,k}=\dim_K\ker p(T)^k\)
jak w
Twierdzeniu~\ref{thm:lin20-kernel-count},
to równoważnie
\[
 \dim_K\mathcal C(T)=
 \sum_p\frac{1}{d_p}
 \sum_{k\ge1}(D_{p,k}-D_{p,k-1})^2.
\]
Przy \(V=\{0\}\) sumy są puste
i mają wartość \(0\).
\end{twierdzenie}
\begin{proof}
Z Twierdzenia~\ref{thm:lin20-elementary}
mamy izomorfizm modułów
\[
 V_T\cong
 \bigoplus_p\bigoplus_{i=1}^{r_p}
           K[t]/(p^{a_{p,i}}).
\]
Tak jak w dowodzie
Twierdzenia~\ref{thm:lin21-invariant-dimension},
wymiar komutanta jest sumą wymiarów
wszystkich przestrzeni homomorfizmów
pomiędzy składnikami tej sumy.

Dla różnych monicznych nierozkładalnych
wielomianów \(p\ne q\)
mamy \(\gcd(p^a,q^b)=1\).
Twierdzenie~\ref{thm:lin21-hom-cyclic}
daje wtedy wymiar homomorfizmów
równy \(0\), a zatem
wszystkie takie składowe znikają.
Dla tego samego \(p\)
zachodzi
\[
 \gcd(p^a,p^b)=p^{\min(a,b)},
 \quad
 \deg\gcd(p^a,p^b)=d_p\min(a,b).
\]
Zsumowanie po \(p\) oraz
po parach bloków daje pierwszy wzór.

Dla dodatnich liczb całkowitych \(a,b\)
liczba indeksów \(k\ge1\)
spełniających jednocześnie \(k\le a\)
i \(k\le b\) jest równa \(\min(a,b)\).
Dlatego suma po wszystkich
uporządkowanych parach \((i,j)\)
jest liczbą trójek \((i,j,k)\)
takich, że
\(1\le k\le a_{p,i}\) i \(k\le a_{p,j}\).
Dla ustalonego \(k\) jest
dokładnie \(B_{p,k}\) możliwych
wyborów \(i\) i tyle samo
niezależnych wyborów \(j\).
Po zsumowaniu po \(k\) otrzymujemy
\[
 \sum_{i,j}\min(a_{p,i},a_{p,j})
 =\sum_{k\ge1}B_{p,k}^2.
\]
Ostatni wzór wynika z
Twierdzenia~\ref{thm:lin20-kernel-count},
które daje
\(D_{p,k}-D_{p,k-1}=d_p B_{p,k}\).
Zatem
\(d_p B_{p,k}^2=
 (D_{p,k}-D_{p,k-1})^2/d_p\),
co kończy dowód.
\end{proof}

\begin{wniosek}[Operator diagonalizowalny]
\label{cor:lin21-diagonalizable}
Niech \(T\) będzie endomorfizmem
\(n\)-wymiarowej przestrzeni nad \(K\),
a jego wielomian charakterystyczny
niech rozkłada się nad \(K\)
na czynniki liniowe.
Jeśli \(T\) jest diagonalizowalny,
a \(m_\lambda=\dim_K\ker(T-\lambda I)\)
oznacza wymiar przestrzeni własnej
dla wartości własnej \(\lambda\),
to
\[
 \dim_K\mathcal C(T)=\sum_\lambda m_\lambda^2.
\]
\end{wniosek}
\begin{proof}
Dla operatora diagonalizowalnego
wszystkie bloki Jordana mają rozmiar \(1\)
na mocy
Twierdzenia~\ref{thm:lin14-diag-criterion}.
Z Wniosku~\ref{cor:lin20-jordan}
jego dzielniki elementarne to
\((t-\lambda)\), występujące
\(m_\lambda\) razy.
Dla \(p=t-\lambda\)
mamy \(d_p=1\),
\(B_{p,1}=m_\lambda\)
oraz \(B_{p,k}=0\) dla \(k\ge2\).
Podstawienie do
Twierdzenia~\ref{thm:lin21-elementary-dimension}
daje podany wzór.
\end{proof}

\section{Centralizator odwracalny}
\label{sec:lin21-units}

\begin{stwierdzenie}[Liczba elementów centralizatora cyklicznego]
\label{prop:lin21-finite-units}
Niech \(K=\mathbb F_q\) będzie
ciałem o \(q\) elementach,
niech \(n\ge1\), a
\(T\in\operatorname{End}_K(V)\)
będzie operatorem cyklicznym,
gdzie \(\dim_K V=n\).
Jeżeli
\(m_T=\prod_{i=1}^{s}p_i^{a_i}\),
gdzie \(p_1,\ldots,p_s\in K[t]\)
są parami różnymi monicznymi
wielomianami nierozkładalnymi
i \(a_i\ge1\), to
\[
 |\mathcal Z(T)|
 =q^n\prod_{i=1}^{s}
       \left(1-q^{-\deg p_i}\right).
\]
\end{stwierdzenie}
\begin{proof}
Na mocy Wniosku~\ref{cor:lin21-cyclic}
grupa \(\mathcal Z(T)\)
odpowiada elementom odwracalnym
pierścienia \(K[t]/(m_T)\).
Wobec chińskiego twierdzenia
o resztach z
Twierdzenia~\ref{thm:lin20-crt}
mamy izomorfizm pierścieni
\[
 K[t]/(m_T)\cong
 \prod_{i=1}^{s}K[t]/(p_i^{a_i}).
\]
Izomorfizm z twierdzenia o resztach
zachowuje mnożenie,
co widać wprost na klasach wielomianów.
Element iloczynu jest odwracalny
dokładnie wtedy, gdy wszystkie
jego współrzędne są odwracalne.

W pierścieniu \(R_i=K[t]/(p_i^{a_i})\)
klasa wielomianu \(h\) jest odwracalna
dokładnie wtedy, gdy \(p_i\nmid h\):
na mocy tożsamości Bézouta
jest to równoważne
\(\gcd(h,p_i^{a_i})=1\).
Pierścień \(R_i\), jako
przestrzeń \(K\)-wektorowa wymiaru
\(a_i\deg p_i\), ma
\(q^{a_i\deg p_i}\) elementów.
Klasy podzielne przez \(p_i\)
tworzą obraz mnożenia przez \(p_i\)
i można je przedstawić jednoznacznie
jako \(p_i h\pmod {p_i^{a_i}}\),
gdzie klasa \(h\) jest liczona
modulo \(p_i^{a_i-1}\).
Jest ich zatem
\(q^{(a_i-1)\deg p_i}\)
(także dla \(a_i=1\), gdy jest
to jedyna klasa zerowa).
Liczba elementów odwracalnych \(R_i\)
jest wobec tego równa
\[
 q^{a_i\deg p_i}
 -q^{(a_i-1)\deg p_i}
 =q^{a_i\deg p_i}
     \left(1-q^{-\deg p_i}\right).
\]
Mnożąc te liczby po \(i\)
i korzystając z
\(\sum_i a_i\deg p_i=\deg m_T=n\),
otrzymujemy tezę.
\end{proof}

\section{Przykłady pełnych obliczeń}
\label{sec:lin21-examples}

\begin{przyklad}[Dwa nilpotentne bloki i pięć parametrów]
\label{prz:lin21-nilpotent}
Niech
\[
 A=J_2(0)\oplus J_1(0)
 =\begin{pmatrix}0&1&0\\0&0&0\\0&0&0\end{pmatrix}.
\]
Dla dowolnej macierzy
\(X=(x_{ij})_{1\le i,j\le3}\)
mamy
\[
 AX=\begin{pmatrix}
 x_{21}&x_{22}&x_{23}\\0&0&0\\0&0&0
 \end{pmatrix},
 \qquad
 XA=\begin{pmatrix}
 0&x_{11}&0\\0&x_{21}&0\\0&x_{31}&0
 \end{pmatrix}.
\]
Warunek \(AX=XA\) jest więc
równoważny
\(x_{21}=x_{23}=x_{31}=0\)
oraz \(x_{22}=x_{11}\).
Każdy element komutanta ma postać
\[
 X=\begin{pmatrix}
 a&b&c\\0&a&0\\0&d&e
 \end{pmatrix},
 \qquad a,b,c,d,e\in K,
\]
i odwrotnie każda taka macierz
przemienia się z \(A\).
Stąd \(\dim_K\mathcal C(A)=5\).
Ten sam wynik otrzymujemy
z dzielników elementarnych \(t^2,t\):
\(B_{t,1}=2\) i \(B_{t,2}=1\),
więc \(2^2+1^2=5\).

Wyznacznik tej macierzy wynosi
\(\det X=a^2e\).
Dlatego \(X\) jest odwracalna
dokładnie wtedy, gdy \(a\ne0\)
i \(e\ne0\).
Jeżeli \(K=\mathbb F_q\),
to \(\mathcal Z(A)\) liczy
\((q-1)^2q^3\) elementów:
parametry \(a,e\) wybieramy
na \(q-1\) sposobów każdy,
a \(b,c,d\) dowolnie.
\end{przyklad}

\begin{przyklad}[Centralizator macierzy diagonalnej]
\label{prz:lin21-diagonal}
Niech \(K=\mathbb Q\)
i \(A=\operatorname{diag}(1,1,2)\).
Z warunku \((\lambda_j-\lambda_i)x_{ij}=0\)
z Przykładu~\ref{prz:lin21-scalar}
wynika
\[
 \mathcal C(A)=
 \left\{\begin{pmatrix}
 a&b&0\\c&d&0\\0&0&e
 \end{pmatrix}:a,b,c,d,e\in\mathbb Q\right\}.
\]
Wobec tego
\(\dim_{\mathbb Q}\mathcal C(A)=5\).
Jego czynniki niezmiennicze to
\(f_1=t-1\) i
\(f_2=(t-1)(t-2)\), jak w
Przykładzie~\ref{prz:lin19-diagonal}.
Wzór z
Twierdzenia~\ref{thm:lin21-invariant-dimension}
daje
\(3\deg f_1+\deg f_2=3+2=5\).
Macierz z tego komutanta
jest odwracalna dokładnie wtedy,
gdy \((ad-bc)e\ne0\).
\end{przyklad}

\begin{przyklad}[Obrót nad ciałem rzeczywistym]
\label{prz:lin21-rotation}
Niech
\[
 A=\begin{pmatrix}0&-1\\1&0\end{pmatrix}
 \in\operatorname{Mat}_{2\times2}(\mathbb R).
\]
Z Przykładu~\ref{prz:lin19-real-rotation}
wiemy, że \(\chi_A=m_A=t^2+1\);
macierz \(A\) jest więc cykliczna.
Z Wniosku~\ref{cor:lin21-cyclic}
otrzymujemy
\[
 \mathcal C(A)=
 \{aI_2+bA:a,b\in\mathbb R\}
 =\left\{\begin{pmatrix}a&-b\\b&a\end{pmatrix}
             :a,b\in\mathbb R\right\}.
\]
Jest to dwuwymiarowa
\(\mathbb R\)-algebra izomorficzna
z \(\mathbb C\) przez
\(a+ib\mapsto aI_2+bA\):
równość \(A^2=-I_2\)
zapewnia zgodność mnożenia.
Wyznacznik wynosi \(a^2+b^2\),
zatem centralizator odwracalny
obejmuje wszystkie takie macierze
poza macierzą zerową.
\end{przyklad}

\begin{przyklad}[Dwa operatory o tych samych wielomianach]
\label{prz:lin21-same-polys}
Niech
\[
 A=J_3(0)\oplus J_3(0),\qquad
 B=J_3(0)\oplus J_2(0)\oplus J_1(0)
\]
nad dowolnym ciałem \(K\).
Jak w Przykładzie~\ref{prz:lin20-not-enough},
obie macierze mają
\(\chi_A=\chi_B=t^6\)
i \(m_A=m_B=t^3\), ale
ich dzielniki elementarne są różne.

Dla \(A\) liczby bloków
o długości co najmniej \(k=1,2,3\)
wynoszą kolejno \(2,2,2\).
Twierdzenie~\ref{thm:lin21-elementary-dimension}
daje więc
\(\dim\mathcal C(A)=2^2+2^2+2^2=12\).
Dla \(B\) odpowiednie liczby
wynoszą \(3,2,1\), skąd
\(\dim\mathcal C(B)=3^2+2^2+1^2=14\).
Wymiary komutantów rozróżniają
te macierze, choć ich wielomiany
charakterystyczne i minimalne są równe.
\end{przyklad}

\begin{przyklad}[Cykliczny blok nad \(\mathbb F_2\)]
\label{prz:lin21-finite}
Niech \(p=t^2+t+1\in\mathbb F_2[t]\)
i \(A=C(p^2)\), jak w
Przykładzie~\ref{prz:lin20-finite}.
Operator jest cykliczny,
ponieważ jego wielomian
minimalny i charakterystyczny
są równe \(p^2\), stopnia \(4\).
Stąd \(\dim_{\mathbb F_2}\mathcal C(A)=4\),
a komutant liczy \(2^4=16\) elementów.

Wielomian \(p\) jest nierozkładalny
nad \(\mathbb F_2\) stopnia \(2\).
Stwierdzenie~\ref{prop:lin21-finite-units}
daje dokładnie
\[
 |\mathcal Z(A)|
 =2^4(1-2^{-2})=12.
\]
Można też policzyć to bezpośrednio:
w pierścieniu \(\mathbb F_2[t]/(p^2)\)
są \(16\) klas wielomianów
stopnia co najwyżej \(3\);
nieodwracalne są dokładnie
cztery klasy podzielne przez \(p\).
\end{przyklad}

\section{Podsumowanie}
\label{sec:lin21-summary}

Komutant operatora \(T\) jest
algebrą \(\operatorname{End}_{K[t]}(V_T)\).
Przy jednym bloku Frobeniusa
jest to dokładnie algebra
wielomianów w \(T\),
izomorficzna z \(K[t]/(m_T)\).
Dla większej liczby bloków
wymiar komutanta otrzymujemy
z przestrzeni homomorfizmów
między modułami cyklicznymi,
których wymiary są równe stopniom
największych wspólnych dzielników.

Dzielniki elementarne dostarczają
równoważnego wzoru
\(\sum_p\deg(p)\sum_k B_{p,k}^2\),
a liczby \(B_{p,k}\) wyliczamy
z różnic wymiarów jąder \(p(T)^k\).
Pozwala to obliczać komutant
bez rozwiązywania pełnego układu
równań \(XT=TX\).
Elementy odwracalne komutanta
tworzą grupę centralizatora.
